% !TeX root = ../../Book1.tex
\section{卷积与光滑化}
\label{b1:sec:2003}

\subsecref{b1:subsec:130403}已经用可积核构造函数的逼近恒等。本节把平均推广到分布：把核的平移看成测试函数，由分布作用于它。所得对象虽然是普通光滑函数，其导数却能保留原分布的微分信息。整个构造先在远离边界的区域进行，不预设分布可以延拓到开集之外。

% 实读 Hunter2001Applied 作者官方 ch11.pdf，11.3--11.4，印刷 pp.297--301。
% 原书讨论 Schwartz 测试函数与缓增分布，11.23 只有稠密性证明概要；
% 本节一般开集 D' 的有限阶差商、共同紧支集 Riemann 和及局部化证明另行完整组织。
% 名称依据：Oleinik2008PDEChinese官方目录1.2用“磨光函数”；只核名称，未读该节定义。
% “磨光子”另见北京大学公开《实变函数》前言，来源记录限定到名称，不据此声称核对其归一化。

\subsection{磨光函数}
\label{b1:subsec:200301}

取非负偶函数 \(\rho\in C_c^\infty(\mathbb R^d)\)，使
\[
 \operatorname{supp}\rho\subset\overline{B(0,1)},
 \quad \int_{\mathbb R^d}\rho=1,
 \quad \rho_\varepsilon(x)\coloneqq\varepsilon^{-d}\cdot \rho(x/\varepsilon).
\]
满足上述条件的核称为本章的\term[moguang-hanshu]{磨光函数}[mollifier]。名称可对照 Oleinik 的《偏微分方程讲义》\cite[第1.2节]{Oleinik2008PDEChinese}；中文资料也使用“磨光子”这一名称。\termidx[moguangzi]{磨光子}\footnote{北京大学数学学院公开的《实变函数》前言，在介绍\chapref{b1:ch:04}的 \(L^p\) 与 Sobolev 空间工具时采用“磨光子”。该用词说明见\href{https://math101.pku.edu.cn/docs/2025-04/b74312fec0b54f0c90ff66270f2bded8.pdf}{学校公布的前言}；这里仅提示名称，不据前言推断其核函数的全部约定。}

\ProseParagraphBreak

例如可将 \(|x|<1\) 上的 \(\exp\bigl(-1/(1-|x|^2)\bigr)\) 在球外补零，再除以其正积分；光滑性已在\secref{b1:sec:2001}的截断构造中说明。这是\cref{b1:def:approximate-identity}中的一种特殊核，光滑性和紧支集是本节新增的要求，偶性则用来简化反射记号。

\ProseParagraphBreak

对开集 \(\Omega\subset\mathbb R^d\)，记
\[
 \Omega_\varepsilon
 \coloneqq\{\,x\in\Omega\mid\operatorname{dist}(x,\Omega^c)>\varepsilon\,\}.
\]
约定到空集的距离为无穷，故 \(\Omega=\mathbb R^d\) 时 \(\Omega_\varepsilon=\mathbb R^d\)。若 \(x\in\Omega_\varepsilon\)，则 \(y\mapsto\rho_\varepsilon(x-y)\) 的支集紧含于 \(\Omega\)；这才保证一般分布可以作用于它。

\subsection{分布卷积}
\label{b1:subsec:200302}

\begin{definition}\label{b1:def:distribution-test-convolution}
设 \(T\in\mathcal D'(\Omega)\)、\(\eta\in C_c^\infty(\mathbb R^d)\)，且 \(\operatorname{supp}\eta\subset\overline{B(0,a)}\)，其中 \(a>0\)。在 \(\Omega_a\) 上定义 \(T\) 与测试函数 \(\eta\) 的\term[fenbu-juanji]{分布卷积}[convolution of a distribution with a test function]为
\[
 (T*\eta)(x)\coloneqq\langle T_y,\eta(x-y)\rangle.
\]
下标 \(y\) 指明分布所作用的变量，不表示对 \(T\) 作逐点取值。
\end{definition}
\symidx{\(T*\eta\)：分布与测试函数的卷积}

若 \(T=T_f\) 来自 \(f\in L^1_{\mathrm{loc}}(\Omega)\)，该定义成为
\[
 (T_f*\eta)(x)=\int_\Omega f(y)\cdot\eta(x-y)\,\mathrm dy.
\]
被积函数支集紧含于 \(\Omega\)，所以每个 \(x\in\Omega_a\) 处都绝对可积，并与前文函数卷积一致。这里没有定义两个任意分布的卷积；即使两个常函数 \(1\) 的通常卷积，在全空间也会发散。紧支集等附加条件可以支持更广的分布卷积，但不参与本节证明。

\begin{proposition}\label{b1:prop:distribution-convolution-smooth}
上述 \(T*\eta\) 属于 \(C^\infty(\Omega_a)\)。对任意多重指标 \(\alpha\)，在 \(\Omega_a\) 上有
\[
 D^\alpha(T*\eta)=T*(D^\alpha\eta)=(D^\alpha T)*\eta.
\]
\end{proposition}
\begin{proof}
固定 \(K\Subset\Omega_a\)，取 \(K\) 的稍大紧邻域 \(K'\Subset\Omega_a\)。当 \(x\in K'\) 时，所有 \(\eta(x-\cdot)\) 的支集都在紧集 \(H=K'+\overline{B(0,a)}\Subset\Omega\) 中。由\eqref{b1:eq:distribution-local-order}，存在 \(C,m\) 使
\[
 |\langle T,\psi\rangle|\le C\cdot p_{H,m}(\psi)
 \quad(\psi\in\mathcal D_H).
\]
对每个坐标方向 \(e_j\)，Taylor 公式及 \(\eta\) 各阶导数的有界性给出
\[
 p_{H,m}\left(
 \frac{\eta(x+he_j-\cdot)-\eta(x-\cdot)}h
 -(\partial_j\eta)(x-\cdot)\right)\le C_\eta|h|
\]
对 \(x\in K\) 和充分小的 \(h\) 一致成立。将 \(T\) 作用于差商，便证明第一阶导数公式；同一估计也证明所得导数连续。把 \(\eta\) 换成其各阶导数，迭代即得全部光滑性与第一项等式。

分布导数的定义又给出
\[
 \langle D^\alpha T_y,\eta(x-y)\rangle
 =(-1)^{|\alpha|}\langle T_y,D_y^\alpha\eta(x-y)\rangle
 =\langle T_y,(D^\alpha\eta)(x-y)\rangle,
\]
因为对 \(y\) 微分还产生一个 \((-1)^{|\alpha|}\)。这证明第二项等式。
\end{proof}

例如在全空间，\(\delta_b*\eta=\eta(\cdot-b)\)，而 \((D^\alpha\delta_b)*\eta=D^\alpha\eta(\cdot-b)\)。点质量及其导数被变成普通光滑函数，但导数阶数没有丢失。

\subsection{正则化}
\label{b1:subsec:200303}

为了讨论收敛，还须把光滑函数 \(T*\eta\) 再与测试函数配对。记反射核 \(\check\eta(z)=\eta(-z)\)。反射不可由卷积记号中擅自省去；当 \(\eta=\rho_\varepsilon\) 时，偶性才使它等于原核。

\begin{proposition}\label{b1:prop:distribution-convolution-pairing}
若 \(\varphi\in\mathcal D(\Omega_a)\)，将其在外部补零，则
\[
 \langle T*\eta,\varphi\rangle
 =\langle T,\check\eta*\varphi\rangle.
\]
右端的卷积属于 \(\mathcal D(\Omega)\)。
\end{proposition}
\begin{proof}
令 \(K=\operatorname{supp}\varphi\)，则
\[
 q(y)=\int_{\mathbb R^d}\varphi(x)\cdot\eta(x-y)\,\mathrm dx
 =(\check\eta*\varphi)(y)
\]
光滑，且支集包含在 \(H=K+\overline{B(0,a)}\Subset\Omega\) 中。以下证明 \(T\) 与这个积分确能交换。

取包含 \(K\) 的有界闭矩形 \(Q\)，并将其分成网格趋于零的小矩形 \(Q_i\)，各选一点 \(x_i\)。函数值 Riemann 和
\[
 q_N(y)=\sum_i m_d(Q_i)\cdot\varphi(x_i)\cdot\eta(x_i-y)
\]
都属于 \(\mathcal D_H\)。对任意固定阶数 \(j\)，\(\varphi(x)D_y^\beta\eta(x-y)\) 在共同紧集上一致连续，\(|\beta|\le j\)，所以 Riemann 和误差关于 \(y\) 一致趋零，即 \(p_{H,j}(q_N-q)\to0\)。有限阶估计因此给出 \(\langle T,q_N\rangle\to\langle T,q\rangle\)。另一方面，由线性性，每个配对都是连续紧支集函数 \(x\mapsto\varphi(x)(T*\eta)(x)\) 的 Riemann 和；在 \(\Omega_a\) 外该乘积补零。它们趋于 \(\int\varphi(x)(T*\eta)(x)\,\mathrm dx\)，等式得证。
\end{proof}

\begin{theorem}[分布的局部光滑化]\label{b1:thm:distribution-regularization}
设 \(T\in\mathcal D'(\Omega)\)，令 \(T_\varepsilon=T*\rho_\varepsilon\in C^\infty(\Omega_\varepsilon)\)。则对每个固定的 \(\varphi\in\mathcal D(\Omega)\)，当 \(\varepsilon\) 充分小时配对有定义，且
\[
 \langle T_\varepsilon,\varphi\rangle\longrightarrow\langle T,\varphi\rangle.
\]
每个 \(D^\alpha T_\varepsilon\) 也以同样的局部意义趋于 \(D^\alpha T\)。
\end{theorem}
\begin{proof}
设 \(K=\operatorname{supp}\varphi\)，取 \(a>0\) 使 \(H=K+\overline{B(0,a)}\Subset\Omega\)。当 \(0<\varepsilon<a\) 时，\(\check\rho_\varepsilon*\varphi\) 和 \(\varphi\) 的支集均在 \(H\) 中，而且
\[
 D^\beta(\check\rho_\varepsilon*\varphi-\varphi)(y)
 =\int\rho(z)\cdot\bigl(D^\beta\varphi(y+\varepsilon z)-D^\beta\varphi(y)\bigr)\,\mathrm dz.
\]
因 \(D^\beta\varphi\) 一致连续、\(\rho\) 只在单位球中有质量，每个阶数的上确界范数都趋于零。前一命题与 \(T\) 的连续性给出所求配对极限。再将 \(T\) 换成 \(D^\alpha T\)，并用卷积微分公式，即得导数结论。
\end{proof}

这里 \(T_\varepsilon\) 的定义域随 \(\varepsilon\) 改变；“局部”指每个固定测试函数最终都能使用，不是把 \(T_\varepsilon\) 未经说明地当成整个 \(\Omega\) 上的函数。也不要求 \(T_\varepsilon\) 本身有逐点极限：例如 \(\delta_0*\rho_\varepsilon=\rho_\varepsilon\)，其峰值通常随 \(\varepsilon^{-d}\) 增长。

\subsection{逼近恒等}
\label{b1:subsec:200304}

\begin{proposition}\label{b1:prop:local-smoothing-approximation}
设 \(K\Subset\Omega\)。若 \(f\in L^p_{\mathrm{loc}}(\Omega)\)、\(1\le p<\infty\)，则
\[
 \|f*\rho_\varepsilon-f\|_{L^p(K)}\longrightarrow0.
\]
若 \(f\in C^k(\Omega)\)，则对所有 \(|\alpha|\le k\)，\(D^\alpha(f*\rho_\varepsilon)\to D^\alpha f\) 在 \(K\) 上一致成立。
\end{proposition}
\begin{proof}
由\secref{b1:sec:2001}的光滑截断，取 \(\chi\in C_c^\infty(\Omega)\)，使它在 \(K\) 的一个邻域上等于 \(1\)。将 \(\chi f\) 补零到全空间；当 \(\varepsilon\) 充分小时，它与 \(f\) 的卷积在 \(K\) 上相同。有限 \(p\) 的结论因此直接归结为\cref{b1:prop:approximate-identity-lp}。

若 \(f\in C^k\)，则 \(\chi f\) 的各阶导数都是连续紧支集函数。对每个 \(|\alpha|\le k\)，卷积微分公式把所需差转为 \((D^\alpha(\chi f))*\rho_\varepsilon-D^\alpha(\chi f)\)，再用同一命题的一致收敛部分即可。
\end{proof}

不能把有限 \(p\) 的结论照搬到 \(p=\infty\)。例如 \(f=\mathbf1_{(0,\infty)}\) 在一维原点附近有跳跃，偶核卷积在原点取值 \(1/2\)，并在附近连续，故其与 \(f\) 的本质上确界距离至少为 \(1/2\)。局部分布收敛、有限 \(p\) 范数收敛和光滑函数的逐阶一致收敛是不同层次，后续使用时须说明所需的是哪一种。

卷积平均怎样消除跳跃并保持局部信息，可见\cref{b1:fig:visual-f15}。
% F15：本书原创矢量示意；依据所在节的定义、证明或明确模型。
\begin{figure}[htbp]
\centering
\begin{tikzpicture}[x=1cm,y=1cm,>=stealth,every node/.style={font=\small},line width=.65pt]
\node[] at (6,3.4) {单位积分光滑核的卷积逼近};
\draw[->,gray] (0.88,0)--(11.8,0) node[right] {\(x\)};\draw[->,gray] (1,-0.12)--(1,2.7);
\draw[gray,dashed] plot coordinates {(3.0000,0.0000) (3.0000,2.2000) (7.0000,2.2000) (7.0000,0.0000)};
\draw[AnalysisBlue,thick] plot coordinates {(0.8000,0.0000) (0.8384,0.0000) (0.8767,0.0000) (0.9151,0.0000) (0.9534,0.0000) (0.9918,0.0000) (1.0301,0.0000) (1.0685,0.0000) (1.1068,0.0000) (1.1452,0.0000) (1.1836,0.0000) (1.2219,0.0000) (1.2603,0.0000) (1.2986,0.0000) (1.3370,0.0000) (1.3753,0.0000) (1.4137,0.0000) (1.4521,0.0000) (1.4904,0.0000) (1.5288,0.0001) (1.5671,0.0005) (1.6055,0.0018) (1.6438,0.0044) (1.6822,0.0087) (1.7205,0.0150) (1.7589,0.0236) (1.7973,0.0344) (1.8356,0.0474) (1.8740,0.0627) (1.9123,0.0801) (1.9507,0.0997) (1.9890,0.1212) (2.0274,0.1446) (2.0658,0.1698) (2.1041,0.1967) (2.1425,0.2251) (2.1808,0.2550) (2.2192,0.2863) (2.2575,0.3189) (2.2959,0.3527) (2.3342,0.3877) (2.3726,0.4236) (2.4110,0.4605) (2.4493,0.4983) (2.4877,0.5369) (2.5260,0.5762) (2.5644,0.6163) (2.6027,0.6569) (2.6411,0.6981) (2.6795,0.7398) (2.7178,0.7819) (2.7562,0.8244) (2.7945,0.8672) (2.8329,0.9103) (2.8712,0.9536) (2.9096,0.9971) (2.9479,1.0407) (2.9863,1.0844) (3.0247,1.1281) (3.0630,1.1718) (3.1014,1.2153) (3.1397,1.2588) (3.1781,1.3020) (3.2164,1.3451) (3.2548,1.3878) (3.2932,1.4302) (3.3315,1.4722) (3.3699,1.5138) (3.4082,1.5548) (3.4466,1.5952) (3.4849,1.6351) (3.5233,1.6742) (3.5616,1.7126) (3.6000,1.7501) (3.6384,1.7868) (3.6767,1.8224) (3.7151,1.8570) (3.7534,1.8905) (3.7918,1.9227) (3.8301,1.9537) (3.8685,1.9832) (3.9068,2.0112) (3.9452,2.0376) (3.9836,2.0623) (4.0219,2.0852) (4.0603,2.1061) (4.0986,2.1251) (4.1370,2.1419) (4.1753,2.1566) (4.2137,2.1690) (4.2521,2.1791) (4.2904,2.1870) (4.3288,2.1927) (4.3671,2.1965) (4.4055,2.1987) (4.4438,2.1997) (4.4822,2.2000) (4.5205,2.2000) (4.5589,2.2000) (4.5973,2.2000) (4.6356,2.2000) (4.6740,2.2000) (4.7123,2.2000) (4.7507,2.2000) (4.7890,2.2000) (4.8274,2.2000) (4.8658,2.2000) (4.9041,2.2000) (4.9425,2.2000) (4.9808,2.2000) (5.0192,2.2000) (5.0575,2.2000) (5.0959,2.2000) (5.1342,2.2000) (5.1726,2.2000) (5.2110,2.2000) (5.2493,2.2000) (5.2877,2.2000) (5.3260,2.2000) (5.3644,2.2000) (5.4027,2.2000) (5.4411,2.2000) (5.4795,2.2000) (5.5178,2.2000) (5.5562,2.1997) (5.5945,2.1987) (5.6329,2.1965) (5.6712,2.1927) (5.7096,2.1870) (5.7479,2.1791) (5.7863,2.1690) (5.8247,2.1566) (5.8630,2.1419) (5.9014,2.1251) (5.9397,2.1061) (5.9781,2.0852) (6.0164,2.0623) (6.0548,2.0376) (6.0932,2.0112) (6.1315,1.9832) (6.1699,1.9537) (6.2082,1.9227) (6.2466,1.8905) (6.2849,1.8570) (6.3233,1.8224) (6.3616,1.7868) (6.4000,1.7501) (6.4384,1.7126) (6.4767,1.6742) (6.5151,1.6351) (6.5534,1.5952) (6.5918,1.5548) (6.6301,1.5138) (6.6685,1.4722) (6.7068,1.4302) (6.7452,1.3878) (6.7836,1.3451) (6.8219,1.3020) (6.8603,1.2588) (6.8986,1.2153) (6.9370,1.1718) (6.9753,1.1281) (7.0137,1.0844) (7.0521,1.0407) (7.0904,0.9971) (7.1288,0.9536) (7.1671,0.9103) (7.2055,0.8672) (7.2438,0.8244) (7.2822,0.7819) (7.3205,0.7398) (7.3589,0.6981) (7.3973,0.6569) (7.4356,0.6163) (7.4740,0.5762) (7.5123,0.5369) (7.5507,0.4983) (7.5890,0.4605) (7.6274,0.4236) (7.6658,0.3877) (7.7041,0.3527) (7.7425,0.3189) (7.7808,0.2863) (7.8192,0.2550) (7.8575,0.2251) (7.8959,0.1967) (7.9342,0.1698) (7.9726,0.1446) (8.0110,0.1212) (8.0493,0.0997) (8.0877,0.0801) (8.1260,0.0627) (8.1644,0.0474) (8.2027,0.0344) (8.2411,0.0236) (8.2795,0.0150) (8.3178,0.0087) (8.3562,0.0044) (8.3945,0.0018) (8.4329,0.0005) (8.4712,0.0001) (8.5096,0.0000) (8.5479,0.0000) (8.5863,-0.0000) (8.6247,0.0000) (8.6630,0.0000) (8.7014,0.0000) (8.7397,0.0000) (8.7781,0.0000) (8.8164,0.0000) (8.8548,0.0000) (8.8932,0.0000) (8.9315,0.0000) (8.9699,0.0000) (9.0082,0.0000) (9.0466,0.0000) (9.0849,0.0000) (9.1233,0.0000) (9.1616,0.0000) (9.2000,0.0000)};
\node[AnalysisBlue] at (10,2.4) {\(\varepsilon=0.8\)};
\draw[orange!80!black,thick] plot coordinates {(0.8000,0.0000) (0.8384,0.0000) (0.8767,0.0000) (0.9151,0.0000) (0.9534,0.0000) (0.9918,0.0000) (1.0301,0.0000) (1.0685,0.0000) (1.1068,0.0000) (1.1452,0.0000) (1.1836,0.0000) (1.2219,0.0000) (1.2603,0.0000) (1.2986,0.0000) (1.3370,0.0000) (1.3753,0.0000) (1.4137,0.0000) (1.4521,0.0000) (1.4904,0.0000) (1.5288,0.0000) (1.5671,0.0000) (1.6055,0.0000) (1.6438,0.0000) (1.6822,0.0000) (1.7205,0.0000) (1.7589,0.0000) (1.7973,0.0000) (1.8356,0.0000) (1.8740,0.0000) (1.9123,0.0000) (1.9507,0.0000) (1.9890,0.0000) (2.0274,0.0000) (2.0658,0.0000) (2.1041,0.0000) (2.1425,0.0000) (2.1808,0.0000) (2.2192,0.0000) (2.2575,0.0000) (2.2959,0.0012) (2.3342,0.0069) (2.3726,0.0203) (2.4110,0.0425) (2.4493,0.0737) (2.4877,0.1133) (2.5260,0.1606) (2.5644,0.2148) (2.6027,0.2750) (2.6411,0.3405) (2.6795,0.4106) (2.7178,0.4847) (2.7562,0.5621) (2.7945,0.6423) (2.8329,0.7248) (2.8712,0.8091) (2.9096,0.8949) (2.9479,0.9816) (2.9863,1.0688) (3.0247,1.1562) (3.0630,1.2433) (3.1014,1.3297) (3.1397,1.4151) (3.1781,1.4990) (3.2164,1.5809) (3.2548,1.6603) (3.2932,1.7368) (3.3315,1.8098) (3.3699,1.8787) (3.4082,1.9428) (3.4466,2.0014) (3.4849,2.0537) (3.5233,2.0989) (3.5616,2.1361) (3.6000,2.1648) (3.6384,2.1844) (3.6767,2.1954) (3.7151,2.1994) (3.7534,2.2000) (3.7918,2.2000) (3.8301,2.2000) (3.8685,2.2000) (3.9068,2.2000) (3.9452,2.2000) (3.9836,2.2000) (4.0219,2.2000) (4.0603,2.2000) (4.0986,2.2000) (4.1370,2.2000) (4.1753,2.2000) (4.2137,2.2000) (4.2521,2.2000) (4.2904,2.2000) (4.3288,2.2000) (4.3671,2.2000) (4.4055,2.2000) (4.4438,2.2000) (4.4822,2.2000) (4.5205,2.2000) (4.5589,2.2000) (4.5973,2.2000) (4.6356,2.2000) (4.6740,2.2000) (4.7123,2.2000) (4.7507,2.2000) (4.7890,2.2000) (4.8274,2.2000) (4.8658,2.2000) (4.9041,2.2000) (4.9425,2.2000) (4.9808,2.2000) (5.0192,2.2000) (5.0575,2.2000) (5.0959,2.2000) (5.1342,2.2000) (5.1726,2.2000) (5.2110,2.2000) (5.2493,2.2000) (5.2877,2.2000) (5.3260,2.2000) (5.3644,2.2000) (5.4027,2.2000) (5.4411,2.2000) (5.4795,2.2000) (5.5178,2.2000) (5.5562,2.2000) (5.5945,2.2000) (5.6329,2.2000) (5.6712,2.2000) (5.7096,2.2000) (5.7479,2.2000) (5.7863,2.2000) (5.8247,2.2000) (5.8630,2.2000) (5.9014,2.2000) (5.9397,2.2000) (5.9781,2.2000) (6.0164,2.2000) (6.0548,2.2000) (6.0932,2.2000) (6.1315,2.2000) (6.1699,2.2000) (6.2082,2.2000) (6.2466,2.2000) (6.2849,2.1994) (6.3233,2.1954) (6.3616,2.1844) (6.4000,2.1648) (6.4384,2.1361) (6.4767,2.0989) (6.5151,2.0537) (6.5534,2.0014) (6.5918,1.9428) (6.6301,1.8787) (6.6685,1.8098) (6.7068,1.7368) (6.7452,1.6603) (6.7836,1.5809) (6.8219,1.4990) (6.8603,1.4151) (6.8986,1.3297) (6.9370,1.2433) (6.9753,1.1562) (7.0137,1.0688) (7.0521,0.9816) (7.0904,0.8949) (7.1288,0.8091) (7.1671,0.7248) (7.2055,0.6423) (7.2438,0.5621) (7.2822,0.4847) (7.3205,0.4106) (7.3589,0.3405) (7.3973,0.2750) (7.4356,0.2148) (7.4740,0.1606) (7.5123,0.1133) (7.5507,0.0737) (7.5890,0.0425) (7.6274,0.0203) (7.6658,0.0069) (7.7041,0.0012) (7.7425,0.0000) (7.7808,0.0000) (7.8192,0.0000) (7.8575,0.0000) (7.8959,0.0000) (7.9342,0.0000) (7.9726,0.0000) (8.0110,0.0000) (8.0493,0.0000) (8.0877,0.0000) (8.1260,0.0000) (8.1644,0.0000) (8.2027,0.0000) (8.2411,0.0000) (8.2795,0.0000) (8.3178,0.0000) (8.3562,0.0000) (8.3945,0.0000) (8.4329,0.0000) (8.4712,0.0000) (8.5096,0.0000) (8.5479,0.0000) (8.5863,0.0000) (8.6247,0.0000) (8.6630,0.0000) (8.7014,0.0000) (8.7397,0.0000) (8.7781,0.0000) (8.8164,0.0000) (8.8548,0.0000) (8.8932,0.0000) (8.9315,0.0000) (8.9699,0.0000) (9.0082,0.0000) (9.0466,0.0000) (9.0849,0.0000) (9.1233,0.0000) (9.1616,0.0000) (9.2000,0.0000)};
\node[orange!80!black] at (10,1.7999999999999998) {\(\varepsilon=0.4\)};
\draw[black,thick] plot coordinates {(0.8000,0.0000) (0.8384,0.0000) (0.8767,0.0000) (0.9151,0.0000) (0.9534,0.0000) (0.9918,0.0000) (1.0301,0.0000) (1.0685,0.0000) (1.1068,0.0000) (1.1452,0.0000) (1.1836,0.0000) (1.2219,0.0000) (1.2603,0.0000) (1.2986,0.0000) (1.3370,0.0000) (1.3753,0.0000) (1.4137,0.0000) (1.4521,0.0000) (1.4904,0.0000) (1.5288,0.0000) (1.5671,0.0000) (1.6055,0.0000) (1.6438,0.0000) (1.6822,0.0000) (1.7205,0.0000) (1.7589,0.0000) (1.7973,0.0000) (1.8356,0.0000) (1.8740,0.0000) (1.9123,0.0000) (1.9507,0.0000) (1.9890,0.0000) (2.0274,0.0000) (2.0658,0.0000) (2.1041,0.0000) (2.1425,0.0000) (2.1808,0.0000) (2.2192,0.0000) (2.2575,0.0000) (2.2959,0.0000) (2.3342,0.0000) (2.3726,0.0000) (2.4110,0.0000) (2.4493,0.0000) (2.4877,0.0000) (2.5260,0.0000) (2.5644,0.0000) (2.6027,0.0000) (2.6411,0.0000) (2.6795,0.0000) (2.7178,0.0000) (2.7562,0.0113) (2.7945,0.0762) (2.8329,0.2000) (2.8712,0.3684) (2.9096,0.5677) (2.9479,0.7869) (2.9863,1.0168) (3.0247,1.2495) (3.0630,1.4772) (3.1014,1.6917) (3.1397,1.8834) (3.1781,2.0406) (3.2164,2.1488) (3.2548,2.1959) (3.2932,2.2000) (3.3315,2.2000) (3.3699,2.2000) (3.4082,2.2000) (3.4466,2.2000) (3.4849,2.2000) (3.5233,2.2000) (3.5616,2.2000) (3.6000,2.2000) (3.6384,2.2000) (3.6767,2.2000) (3.7151,2.2000) (3.7534,2.2000) (3.7918,2.2000) (3.8301,2.2000) (3.8685,2.2000) (3.9068,2.2000) (3.9452,2.2000) (3.9836,2.2000) (4.0219,2.2000) (4.0603,2.2000) (4.0986,2.2000) (4.1370,2.2000) (4.1753,2.2000) (4.2137,2.2000) (4.2521,2.2000) (4.2904,2.2000) (4.3288,2.2000) (4.3671,2.2000) (4.4055,2.2000) (4.4438,2.2000) (4.4822,2.2000) (4.5205,2.2000) (4.5589,2.2000) (4.5973,2.2000) (4.6356,2.2000) (4.6740,2.2000) (4.7123,2.2000) (4.7507,2.2000) (4.7890,2.2000) (4.8274,2.2000) (4.8658,2.2000) (4.9041,2.2000) (4.9425,2.2000) (4.9808,2.2000) (5.0192,2.2000) (5.0575,2.2000) (5.0959,2.2000) (5.1342,2.2000) (5.1726,2.2000) (5.2110,2.2000) (5.2493,2.2000) (5.2877,2.2000) (5.3260,2.2000) (5.3644,2.2000) (5.4027,2.2000) (5.4411,2.2000) (5.4795,2.2000) (5.5178,2.2000) (5.5562,2.2000) (5.5945,2.2000) (5.6329,2.2000) (5.6712,2.2000) (5.7096,2.2000) (5.7479,2.2000) (5.7863,2.2000) (5.8247,2.2000) (5.8630,2.2000) (5.9014,2.2000) (5.9397,2.2000) (5.9781,2.2000) (6.0164,2.2000) (6.0548,2.2000) (6.0932,2.2000) (6.1315,2.2000) (6.1699,2.2000) (6.2082,2.2000) (6.2466,2.2000) (6.2849,2.2000) (6.3233,2.2000) (6.3616,2.2000) (6.4000,2.2000) (6.4384,2.2000) (6.4767,2.2000) (6.5151,2.2000) (6.5534,2.2000) (6.5918,2.2000) (6.6301,2.2000) (6.6685,2.2000) (6.7068,2.2000) (6.7452,2.1959) (6.7836,2.1488) (6.8219,2.0406) (6.8603,1.8834) (6.8986,1.6917) (6.9370,1.4772) (6.9753,1.2495) (7.0137,1.0168) (7.0521,0.7869) (7.0904,0.5677) (7.1288,0.3684) (7.1671,0.2000) (7.2055,0.0762) (7.2438,0.0113) (7.2822,0.0000) (7.3205,0.0000) (7.3589,0.0000) (7.3973,0.0000) (7.4356,0.0000) (7.4740,0.0000) (7.5123,0.0000) (7.5507,0.0000) (7.5890,0.0000) (7.6274,0.0000) (7.6658,0.0000) (7.7041,0.0000) (7.7425,0.0000) (7.7808,0.0000) (7.8192,0.0000) (7.8575,0.0000) (7.8959,0.0000) (7.9342,0.0000) (7.9726,0.0000) (8.0110,0.0000) (8.0493,0.0000) (8.0877,0.0000) (8.1260,0.0000) (8.1644,0.0000) (8.2027,0.0000) (8.2411,0.0000) (8.2795,0.0000) (8.3178,0.0000) (8.3562,0.0000) (8.3945,0.0000) (8.4329,0.0000) (8.4712,0.0000) (8.5096,0.0000) (8.5479,0.0000) (8.5863,0.0000) (8.6247,0.0000) (8.6630,0.0000) (8.7014,0.0000) (8.7397,0.0000) (8.7781,0.0000) (8.8164,0.0000) (8.8548,0.0000) (8.8932,0.0000) (8.9315,0.0000) (8.9699,0.0000) (9.0082,0.0000) (9.0466,0.0000) (9.0849,0.0000) (9.1233,0.0000) (9.1616,0.0000) (9.2000,0.0000)};
\node[black] at (10,1.2) {\(\varepsilon=0.15\)};

\end{tikzpicture}
\caption[卷积如何移动、平均与光滑化]{卷积如何移动、平均与光滑化。虚线为区间示性函数；实线为它与尺度化光滑核的卷积。核取 \(C\exp(-1/(1-x^2))\) 于 \(|x|<1\)，区间外为零，其中常数使积分为一。}
\label{b1:fig:visual-f15}
\end{figure}

